\section{Conclusion}\label{sec:conclusion}

HHL is verified in Isabelle/HOL from the spectral theorem to a list of two-qubit gates: phase
estimation with its precision, the exact algorithm as a matrix identity and as a totally correct
program, the paper's error analysis with explicit constants, and, for the shifted cycle Laplacian, a
complete gate-level pipeline whose output is provably close to the normalised solution and whose length
is a closed-form polynomial. The formalisation corrected seven statements of the original paper.

The development was written by an agent and directed by a person. The division of labour that
emerged is sharp. The agent wrote the proofs, refactored without fatigue, and resumed after five months
by reading what had been written down; the human chose what to prove, checked that the statements meant
what the paper meant, decided names, structure and presentation, and insisted on the checks that caught
the agent's mistakes. The kernel guarantees that the theorems are true; it is the human who
guarantees that they are the right theorems.

\paragraph{Future work.} The entries for the lectures and HHL await the review of the functional
calculus at the AFP. Beyond them lie the later quantum linear-system algorithms~\cite{childs2017}, gate
counts for other families of matrices, and a bridge that removes the remaining distinction between the
matrix world of the lectures and the operator world of Unruh's libraries.

\paragraph{Use of AI.} As described in Section~\ref{sec:process}, the formal development was written
by AI agents (Claude Code with Claude Opus~4.7, Claude Fable~5.1 and Claude Opus~5.5; OpenAI Codex with
GPT-5.4) under the author's direction. Claude Code also measured the statistics of Section~\ref{sec:process}
and drafted this paper with the author; every number in it is computed from the repository and the
logs as described in Appendix~\ref{app:method}.
